We project pathway features to EEG-like signals via a learned linear transformation:

\begin{equation}
\mathcal{E} = \mathbf{W} \mathcal{P} + \mathbf{b} + \boldsymbol{\epsilon}
\end{equation}

where $\mathbf{W} \in \mathbb{R}^{105 \times 6}$ is a learned projection matrix, $\mathbf{b} \in \mathbb{R}^{105}$ is a bias vector, and $\boldsymbol{\epsilon} \sim \mathcal{N}(0, \sigma^2)$ is Gaussian noise ($\sigma = 0.1$).

\subsection*{Design Rationale}

\begin{enumerate}
    \item \textbf{Linearity}: Real EEG involves complex filtering and volume conduction, but the leading-order effect is linear spatial mixing. Our projection models this directly.
    
    \item \textbf{105 Channels}: Standard EEG systems use $\sim 64$--$128$ channels. We fix 105 to match the 10-5 electrode system used in clinical EEG.
    
    \item \textbf{Noise Injection}: Real recordings are corrupted by thermal noise, muscle artifact, and other sources. We add Gaussian noise with $\sigma = 0.1$ (10\% of typical signal magnitude).
    
    \item \textbf{No Temporal Dynamics}: We project static pathway signatures, not time series. This is analagous to extracting EEG features at a single time point.
\end{enumerate}

\subsection*{Learning Objective}

The projection matrix is learned jointly with predictors via standard regression:

\begin{equation}
\mathcal{L}_{\text{proj}} = \text{MSE}(\hat{s}_{\mathcal{P}}, s) + \text{MSE}(\hat{s}_{\mathcal{E}}, s)
\end{equation}

This encourages the projector to preserve task-relevant information about sparsity.
